\section*{Hoofdstuk \ref{ch:b2:oceanic-islands} --- Oceaaneilanden}

\begin{solution}{exo:b2:oceanic-islands:1}
(a) Preschildstadium, zoals Loihi vandaag. (b) \omterm{def:b2:oceanic-islands:stages}{Schildstadium}. (c) Postschildstadium. (d) Verjongd
stadium. (e) Geen eruptief stadium meer: de vulkaan is onder zijn rif gezonken, dat nu een \omterm{def:b2:oceanic-islands:reefs}{atol} is.
\end{solution}

\begin{solution}{exo:b2:oceanic-islands:2}
(a) $2500 + 350\sqrt9 = \qty{3550}{m}$. (b) $350\times(\sqrt{25} - \sqrt9) = 350\times2 =
\qty{700}{m}$. (c) $350/(2\times5) = \qty{35}{m}$ per miljoen jaar.
\end{solution}

\begin{solution}{exo:b2:oceanic-islands:3}
(a) $5000/17 = \qty{294}{kyr}$ rond de hoofdeilanden; $30\,000/70 = \qty{429}{kyr}$ langs de
rug. (b) Er worden afschuivingen gemist: een jongere vulkaan kan het litteken en het puin van een oudere begraven,
en oude afzettingen worden door sediment bedekt. Het werkelijke aantal is groter en het werkelijke interval
korter.
\end{solution}

\begin{solution}{exo:b2:oceanic-islands:4}
(a) $4158\times0.3048 = \qty{1267}{m}$; $4610\times0.3048 = \qty{1405}{m}$; $150\times0.3048 =
\qty{46}{m}$. (b) Het eerste rif groeide niet dieper dan \qty{46}{m} en zijn basis ligt nu op
\qty{1267}{m}: de vulkaan is minstens $1267 - 46 = \qty{1221}{m}$ gezonken (meer als de zeespiegel
toen hoger stond). (c) Het wegzinken is hoogstens ongeveer \qty{1267}{m}, over minstens \qty{33.9}{Myr}:
hoogstens $1267/33.9 = \qty{37}{m}$ per miljoen jaar, \qty{0.037}{mm} per jaar.
\end{solution}

\begin{solution}{exo:b2:oceanic-islands:5}
(a) $\dot S = 175/\sqrt t$: \qty{87.5}{m.Myr^{-1}} (\qty{0.0875}{mm.yr^{-1}}) bij \qty{4}{Myr},
\qty{35}{m.Myr^{-1}} (\qty{0.035}{mm.yr^{-1}}) bij \qty{25}{Myr}, \qty{25}{m.Myr^{-1}}
(\qty{0.025}{mm.yr^{-1}}) bij \qty{49}{Myr}. (b) $4.2/0.0875 = 48$ keer. (c) $2.6/0.0875 = 30$
keer: de belastingsfase laat een eiland veel sneller zinken dan de afkoeling van een jonge plaat.
\end{solution}

\begin{solution}{exo:b2:oceanic-islands:6}
(a) $\dot S = (3200/62.8)\,e^{-100/62.8} = 50.96\times0.2035 = \qty{10.37}{m}$ per miljoen jaar,
\qty{0.0104}{m} per duizend jaar: de gemeten \qty{0.01}{m} per duizend jaar. (b) $e^{-t/62.8}
= 10/50.96$, $t = 62.8\ln5.096 = \qty{102.3}{Myr}$. (c) Wortelwet: $2500 + 350\times5 =
\qty{4250}{m}$; exponentiële wet: $6400 - 3200\times0.6716 = \qty{4251}{m}$; ze verschillen
\qty{1}{m}.
\end{solution}

\begin{solution}{exo:b2:oceanic-islands:7}
(a) $\qty{8.6}{cm.yr^{-1}} = \qty{86}{km.Myr^{-1}}$; $86\times33 = \qty{2838}{km}$. (b) Ja:
\qty{27.7}{Myr} is minder dan de \qty{33}{Myr} waarna de riffen verdrinken, dus zijn rif zou
nog moeten leven. (c) $55\times1.609 = \qty{88.5}{km}$, of $88.5/86 = \qty{1.0}{Myr}$ meer: ongeveer
\qty{28.7}{Myr}. Dat is iets minder dan \qty{33}{Myr}, maar de \qty{33}{Myr} is zelf een
benaderend getal voor de hele keten, en de groei hangt ook af van de watertemperatuur bij elk
\omterm{def:b2:oceanic-islands:reefs}{atol}: Kure, het laatste levende \omterm{def:b2:oceanic-islands:reefs}{atol}, ligt inderdaad dicht bij het \omterm{def:b2:oceanic-islands:darwin-point}{Darwinpunt}.
\end{solution}

\begin{solution}{exo:b2:oceanic-islands:8}
(a) Doos: $30\times17\times2 = \qty{1020}{km^3}$; piramide: een derde, \qty{340}{km^3}. (b) Eén blok
alleen is een tiende (doos) of een dertigste (piramide) van de hele vulkaan Kilauea. (c) $23\,000\times
0.1 = \qty{2300}{km^3}$, een kwart van de Kilauea; een schatting die staat of valt met de aangenomen
dikte.
\end{solution}

\begin{solution}{exo:b2:oceanic-islands:9}
(a) $1000/2.6 = \qty{385}{kyr}$. (b) Ja: \qty{2.6}{m} per duizend jaar is
\qty{2.6}{mm.yr^{-1}}, minder dan \qty{4.2}{mm.yr^{-1}}, een marge van \qty{1.6}{mm.yr^{-1}}.
(c) Relatieve stijging $12 + 2.6 = \qty{14.6}{mm.yr^{-1}}$; het rif raakt $14.6 - 4.2 =
\qty{10.4}{mm.yr^{-1}}$ achter gedurende $16.5 - 8.2 = \qty{8.3}{kyr}$: zijn top eindigt op $10.4\times8.3 =
\qty{86}{m}$. (d) $(12 - 4.2)\times8.3 = \qty{65}{m}$. In beide gevallen blijft het boven de diepste
groeigrens, maar in veel donkerder water: of het weer bijhoudt, hangt ervan af hoe zijn groei met de
diepte vertraagt.
\end{solution}

\begin{solution}{exo:b2:oceanic-islands:10}
(a) $S(t_1) = C(\sqrt{t_1 + \Delta} - \sqrt{t_1})$ heeft als afgeleide $\frac C2(1/\sqrt{t_1 +
\Delta} - 1/\sqrt{t_1}) < 0$: ze neemt af met $t_1$. (b) $350\sqrt{20} = \qty{1565}{m}$; $350\times
(\sqrt{70} - \sqrt{50}) = 350\times1.2955 = \qty{453}{m}$. (c) Nee: $\sqrt t$ is eindig in $0$, dus
het wegzinken over elk interval is eindig, \qty{1565}{m} in (b). Alleen de snelheid op het ogenblik $t = 0$
is oneindig, waar de wet niet meer bedoeld is om de bodem in detail te beschrijven.
\end{solution}

\begin{solution}{exo:b2:oceanic-islands:11}
(a) $30\times0.1 = \qty{3}{m}$. (b) Ongeveer \qty{120}{m} als zijn top bij de hoogstand op zeeniveau
stond: het \omterm{def:b2:oceanic-islands:reefs}{atol} wordt een kalksteeneiland, regen lost het aragoniet van de koralen op en het
gesteente herkristalliseert tot calciet, met grotten. (c) $120/3 = 40$ cycli, \qty{4}{Myr}. (d) De uitgeloogde
zones zijn deze perioden van blootstelling, verzegeld onder het rif dat opnieuw groeide toen de zee terugkwam.
Over vele cycli keert de zee terug naar vergelijkbare hoogstanden, dus $\Delta\eta$ is klein, de riftop
groeit bij elke hoogstand weer tot zeeniveau aan, en de dikte van de kap wordt bepaald door het wegzinken, zoals de
propositie zegt.
\end{solution}

\begin{solution}{exo:b2:oceanic-islands:12}
(a) $3200\times(e^{-50/62.8} - e^{-120/62.8}) = 3200\times(0.4510 - 0.1480) = \qty{970}{m}$.
(b) $150 + 970 = \qty{1120}{m}$. (c) $700\times1.8288 = \qty{1280}{m}$, \qty{160}{m} dieper dan de
schatting. (d) De \omterm{def:b2:oceanic-islands:seamount}{guyot} kan eerder zijn verdronken, op een jongere plaat, en meer zijn gezonken; een naburige
vulkaan kan hem hebben belast; zijn top kan door golven lager zijn afgeschaafd, of door sediment bedekt zijn; de
plaat eronder kan een andere ouderdom hebben; en het zeeniveau van het tijdstip van verdrinken was niet
dat van vandaag.
\end{solution}

\begin{solution}{pb:b2:oceanic-islands:1}
\textbf{1.} $2500 + 350\sqrt{10} = 2500 + 350\times3.1623 = \qty{3607}{m}$.
\textbf{2.} $3607 + 1000 = \qty{4607}{m}$.
\textbf{3.} $V = \pi r^2h/3 = \pi\times50^2\times4.607/3 = \qty{12061}{km^3}$.
\textbf{4.} $0.85\times12\,061 = \qty{10252}{km^3}$ in \qty{1}{Myr}: \qty{0.0103}{km^3} per jaar,
tien keer minder dan de Kilauea vandaag.
\textbf{5.} $95\,600/12\,061 = 7.9$: ongeveer acht.
\textbf{6.} Tholeïtische lava, uit het \omterm{def:b2:oceanic-islands:stages}{schildstadium}; de alkalische lava's van het postschildstadium dekken
hem af.
\textbf{7.} $\qty{1}{km}/\qty{1}{Myr} = \qty{1}{m}$ per duizend jaar, 2.6 keer trager dan voor
het eiland Hawaï.
\textbf{8.} Ja: $4.2 - 1.0 = \qty{3.2}{mm.yr^{-1}}$ over.
\textbf{9.} $12 + 1 - 4.2 = \qty{8.8}{mm}$ per jaar.
\textbf{10.} $150/8.8 = \qty{17}{kyr}$. Een verdronken rifterras op de flank, een trede van het soort
dat voor het eiland Hawaï is gevonden.
\textbf{11.} $350\times(\sqrt{22} - \sqrt{10}) = 350\times(4.6904 - 3.1623) = \qty{535}{m}$.
\textbf{12.} $1000 - 535 = \qty{465}{m}$, met $465/12 = \qty{38.8}{m}$ per miljoen jaar.
\textbf{13.} $350/(2\sqrt{10}) = \qty{55.3}{m.Myr^{-1}}$ in het begin; $350/(2\sqrt{22}) =
\qty{37.3}{m.Myr^{-1}}$ aan het einde.
\textbf{14.} $2500 + 350\times4.6904 = \qty{4142}{m}$.
\textbf{15.} De grote Hawaï-eilanden rijzen tot ongeveer vier keer hoger op, dus er is veel meer
gesteente boven zeeniveau weg te nemen; hun \omterm{def:b2:oceanic-islands:stages}{schildstadium} duurde ook langer.
\textbf{16.} Golven, die de kust terugdringen en een abrasieplatform afschaven; aardverschuivingen en
\omterm{def:b2:oceanic-islands:collapse}{flankinstortingen}.
\textbf{17.} Een \omterm{def:b2:oceanic-islands:reefs}{barrièrerif}, met een lagune tussen het rif en het laatste stuk vulkanische kust.
\textbf{18.} $350\times(\sqrt{70} - \sqrt{22}) = 350\times(8.3666 - 4.6904) = \qty{1287}{m}$.
\textbf{19.} $1287/48 = \qty{26.8}{m}$ per miljoen jaar, \qty{0.027}{mm} per jaar: $4.2/0.0268 =
157$ keer minder dan het rif kon groeien.
\textbf{20.} \qty{20}{m} dikker dan de kalksteen boven het basalt onder Parry Island
(\qty{1267}{m}), en \qty{118}{m} dunner dan onder Elugelab (\qty{1405}{m}): dezelfde orde van grootte.
\textbf{21.} $2500 + 350\times8.3666 = \qty{5428}{m}$.
\textbf{22.} $3200\times(e^{-70/62.8} - e^{-120/62.8}) = 3200\times(0.3280 - 0.1480) =
\qty{576}{m}$.
\textbf{23.} $150 + 576 = \qty{726}{m}$.
\textbf{24.} De rifkap van het \omterm{def:b2:oceanic-islands:reefs}{atol} Petrel is \textbf{\qty{1287}{m}} dik, ongeveer \qty{1.3}{km}
rifkalksteen, zoals onder Enewetak.
\end{solution}
